% Post Reditum in Senatu (post-reditum-in-senatu) --- English translation
% Translated by Claude (Anthropic) from the Latin (Perseus).
% Style guide: STYLE.md.

\ciceroSection{1} Senators, if for your immortal services to me and to my brother and to our children I shall have given you thanks too little fully, I beg and entreat you to think it is to be set down to my own nature rather than to the size of your kindnesses. For what such fertility of mind can there be, what such abundance of speaking, what such divine and incredible kind of utterance, by which any man could --- I shall not say embrace by speaking, but even survey by counting --- the whole of what you have together earned by giving to us? You have given me back my brother whom I longed for; me to my brother who loves me so much; their parents to our children, our children to us; you have given back rank, order, fortunes, this great commonwealth, this country than which nothing can be more sweet --- and at last, you have given us back to ourselves.

\ciceroSection{2} If we ought to hold our parents most dear, because life and patrimony and freedom and citizenship are handed to us from them; if the immortal gods, by whose kindness we have held these things and been increased in the rest; if the Roman people, by whose offices we have been placed in the greatest council and in the highest grade of standing and in this citadel of all the lands; if this very order, by whose magnificent decrees we have so often been honoured: then it is something measureless and unbounded that we owe to you --- you who, by your singular zeal and singular agreement, have at one and the same moment given us back the kindnesses of our parents, the gifts of the immortal gods, the offices of the Roman people, and your own many judgements about me. So that, while we owe much to you, much to the Roman people, things innumerable to our parents, all things to the immortal gods, those things separately we held before through them; now, all together, we have recovered through you.

\ciceroSection{3} And so, senators, I appear to have obtained through you what is not even to be wished for in a man --- a kind of immortality. For what time will there ever be when the memory and the fame of your kindnesses to me shall die? You who, at the very moment when, hemmed in by violence, by the sword, by fear, by threats, you were held under siege --- not long after my departure recalled me by common voice, on the motion of L. Ninnius, the bravest and best of men, whom that pestilential year had as the most loyal and the least fearful defender of my safety, had it been thought fit to fight. After the power of decreeing was no longer left to you, by the act of that tribune of the plebs who, when he could not in his own person tear at the commonwealth, had hidden himself under another man's villainy --- you were never silent about me, never failed to demand my safety from those consuls who had sold it.

\ciceroSection{4} And so by your zeal and authority it was brought about that that very year which I had wished to be fatal to me rather than to my country had eight tribunes of the plebs who promulgated for my safety and again and again referred the matter to you. For the consuls, modest men and law-abiding, were prevented by a law --- not the law against me, but the law passed against themselves, which my enemy promulgated, that, if those who had almost destroyed our state should rise from the dead, then I should return: by which act he confessed both that he longed for the lives of those men, and that the commonwealth would be in great danger if, with the enemies and assassins of the commonwealth alive again, I had not come back. And in that very year, when I had withdrawn, when the chief man of the citizenry was guarding his life not by the protection of laws but by his own walls, when the commonwealth was without consuls, robbed not only of its perpetual parents but even of its yearly guardians, when you were forbidden to give your opinions, when the heading of my proscription was being read out --- you never doubted to bind my safety together with the common safety.

\ciceroSection{5} But after, by the singular and outstanding courage of the consul P. Lentulus, you began on the Kalends of January to look out from the fog and darkness of the previous year into the light of the commonwealth; when the great standing of Q. Metellus, that most noble man and best of men, had come to the aid of the state, when the courage and good faith of the praetors and almost all the tribunes of the plebs had come to the aid of the state; when Cn. Pompeius --- in courage, in glory, in deeds easily the first of all peoples, of all ages, of all memory --- thought he could come safely into the Senate: such was your common consent about my safety that, while my body was absent, my standing had already returned to my country.

\ciceroSection{6} In that month indeed you were able to judge what difference there was between me and my enemies. I gave up my own safety, that the commonwealth should not be drenched in citizens' wounds on my account; they thought my return must be cut off, not by the votes of the Roman people, but by a flood of blood. And so afterwards you gave no answer to citizens, none to allies, none to kings; nothing the courts declared by their verdicts, nothing the people by their votes, nothing this order by its authority; you saw the Forum mute, the Curia tongueless, the city silent and broken.

\ciceroSection{7} At that very time, when the man had withdrawn who, with you for his backers, had stood up against slaughter and burning --- when you saw men flying about the whole city with iron and torches, the houses of the magistrates assaulted, the temples of the gods set on fire, the rods of fasces of the highest man and most distinguished consul broken, the most sacred body of a most brave and most upstanding tribune of the plebs not merely touched and violated by hand but wounded with iron and cut to pieces --- some magistrates, shaken by that carnage, drew back a little from my cause: some out of fear of death, some out of despair of the commonwealth. The rest there were whom no terror, no force, no hope, no fear, no promises, no threats, no weapons, no torches could drive away from your authority, from the dignity of the Roman people, and from my safety.

\ciceroSection{8} The chief man, P. Lentulus --- parent and god of our life, fortune, memory, and name --- thought this would be the proof of his courage, this the token of his spirit, this the very light of his consulship, if he should give me back to myself, to my own people, to you, and to the commonwealth. As soon as he was named consul-elect, he never doubted to deliver an opinion worthy of my safety and of the commonwealth. When he was forbidden by a tribune of the plebs, when that splendid clause of the law was being read out --- that no man should refer the matter to you, no man decree, no man debate, no man speak, no man cross to vote, no man be present at the writing --- he held that whole proscription to be no law: by which a citizen with the highest deserts of the commonwealth was being torn away by name, without trial, together with the Senate, from the commonwealth itself. And when he entered upon his magistracy, I shall not say what he did first, but what at all he did, except that, by saving me, he should set up your authority and standing for the time to come.

\ciceroSection{9} O immortal gods, what a kindness you seem to me to have given, that this year P. Lentulus was consul of the Roman people! And how much greater you would have given, had he been consul the year before. For I should not have needed the medicine of a consul, had I not fallen by a consul's wound. I had heard from the wisest of men and the best citizen and man, Q. Catulus, that there had often not been a single bad consul, but never two --- with the exception of that Cinnan time. So he used to say that my cause would always be the firmest, while there was even one consul left in the state. He spoke truly, if that condition --- the two-consul case, which had not before been seen in the commonwealth --- could have lasted as a continuing thing. For if Q. Metellus had been consul at that time as my enemy, do you doubt with what spirit he was likely to be in saving me, when you see he has been the prime mover and signatory in restoring me?

\ciceroSection{10} But there were those consuls whose minds, narrow, low, twisted, filled up with darkness and filth, could neither look up to nor sustain nor take in the very name of consul, the splendour of that office, the greatness of so great a command --- not consuls, but traffickers in provinces and sellers of your standing. Of whom one was demanding back from me Catiline --- his lover --- in many men's hearing; the other Cethegus, his cousin. These two, the most villainous since the memory of men, not consuls but bandits, did not merely desert me in a public and consular cause, but betrayed me, attacked me, and wished me to be stripped of every help, not only their own, but yours and that of all the orders. Yet one of them deceived neither me nor anyone.

\ciceroSection{11} For who could entertain any hope of any good in a man whose youth had been openly available to all kinds of lust? Who could not even keep men's filthy unrestraint away from the most sacred parts of his own body? Who, when he had got through his own estate no less briskly than later he got through the commonwealth, kept up his poverty and luxury by playing the pander in his own household? Who, had he not taken refuge on the altar of the tribunate, could not have escaped the praetor's force, his creditors' multitude, or the proscription of his goods? Who, in his magistracy, had he not carried the bill on the pirate war, would surely, driven by poverty and corruption, have turned pirate himself --- and would have done less harm to the commonwealth than that an unspeakable enemy and brigand walked within our walls? It was while he sat watching that a tribune of the plebs carried a law that the auspices should not be obeyed, that no man should announce against an assembly or against a comitia, that no man should veto a law --- that the Aelian and Fufian laws, which our ancestors had wished to be the surest protections of the commonwealth against tribunician madnesses, should have no force.

\ciceroSection{12} And the same man afterwards, when an innumerable multitude of good men in mourning had come down from the Capitol as suppliants to him, when the most noble youths and all the Roman knights had thrown themselves at the feet of the most shameless of pimps --- with what a face the curled-haired libertine spurned not only the tears of the citizens but the prayers of the country itself! And he was not content with that, but mounted to a public meeting and there said what, had Catiline her former husband risen from the dead, he would not have dared to say: that he would exact from the Roman knights the penalties for the Nones of December which had taken place in my consulship and for the Capitoline slope. Nor only did he say this, but he laid hands by name on whom he pleased, and as for L. Lamia, a Roman knight, a man of outstanding standing, most loving of my safety in his friendship, most loving of the commonwealth in his fortunes --- he, the masterful consul, ordered out of the city. And when you had voted that the senatorial dress should be changed, and you all had changed it, and all good men had done so before --- he, smeared with unguents, in the praetexta which all the praetors and aediles had then thrown off, mocked the squalor of you and the mourning of a most dear citizenry; and he did what no tyrant ever did --- said nothing to keep you from groaning over your evil in secret, but issued an edict that you should not openly mourn the country's ills.

\ciceroSection{13} But when in the Circus Flaminius he was led out into the public meeting --- not by a tribune of the plebs as a consul, but by an arch-pirate as a bandit --- he came forward first, with what consular weight! full of wine, of sleep, of debauchery, his hair dripping, his locks dressed, his eyes heavy, his cheeks flopping; in a thick, drunken voice he, as a man of grave authority, declared that the punishment of citizens not condemned by trial displeased him most violently. Where had this great authority of his lain hidden so long? Why had so outstanding a courage in this curled-haired dancing master been on holiday so long, in dens and gluttonies? For that other one, Caesoninus Calventius, has from his youth haunted the Forum, with nothing to recommend him except a feigned and crafty sourness --- no judgement, no abundance of speaking, no zeal for soldiering or for getting to know men, no liberality. Whom, going past, when you had seen him uncombed, bristly, mournful, even if you thought him rustic and uncivil, you would still not have judged him debauched and ruined.

\ciceroSection{14} Whether you stood in the Forum with this man or with a stump, you would have thought there was no difference --- without sense, without taste, tongueless, slow, an inhuman business; you would have called him only a Cappadocian just dragged out of a slave-pen. The same man at home, how lustful, how filthy, how unrestrained --- with pleasures admitted not by the front door but smuggled in by the back! Yet when even he begins to study letters, and the monstrous beast to philosophise with little Greeks, then he is an Epicurean: not deeply given to that doctrine, whatever it is, but caught by the one word ``pleasure.'' And he has masters, not from those impertinent ones who debate whole days about duty and about virtue, who exhort one to labour, to industry, to undergoing dangers for one's country --- but those who argue that no hour ought to be empty of pleasure, that in every part of the body some joy and delight should always be found.

\ciceroSection{15} These he uses as masters of the household for his lusts; these track and sniff out every pleasure; these are the founders and arrangers of the dinner-party; the same men weigh and assess the pleasures, give an opinion and judge how much should be assigned to each lust. Trained in their arts, he so despised this most prudent state that he reckoned all his lusts and all his shamefulness could be hidden, if only he carried his hard face into the Forum. He, in fact, did not deceive me at all --- for I had known, by the marriage tie of the Pisones, how far his maternal birth, of Transalpine blood, had drawn him from this race of men --- but he deceived you and the Roman people, not by judgement nor eloquence (which often happens with many), but by his wrinkles and his eyebrow.

\ciceroSection{16} L. Piso, did you dare with that eye --- I shall not say with that mind --- with that face, not life, with such an eyebrow, for I cannot say with such great deeds done, to combine your counsels with A. Gabinius for my destruction? Did not the smell of his unguents, the panting of his wine, the brow stamped with the tracks of the curling-iron lead you to this thought --- that, when you had been similar to him in fact, you should not be allowed any longer to use the cover of your own brow to hide such great shamefulness? With this man you dared to enter into a compact, that you should sell off, by the bargain of provinces, the consular dignity, the standing of the commonwealth, the authority of the Senate, and the fortunes of a citizen of the highest deserts? In your consulship, by your edicts and orders, the Senate of the Roman people was not allowed to come to the aid of the commonwealth --- not only by their opinions and their authority, but not even by their grief and their dress?

\ciceroSection{17} Did you suppose yourself consul at Capua --- in which city was once the home of pride --- as you were at that time, or at Rome, in whose state all the consuls before you obeyed the Senate? You dared, when led out into the Circus Flaminius with that companion of yours, to say that you had always been merciful? By which word you were demonstrating that the Senate, and all good men, when they were warding the plague off the country, had been cruel. Were you, the merciful, who had set me, your kinsman by marriage, as the foreman of the prerogative century in your election, who had asked me on the Kalends of January to give my opinion in the third place --- were you the merciful when you handed me over, bound, to the enemies of the commonwealth? You spurned my son-in-law, your cousin, with the proudest and cruellest words from your knees; you spurned your kinswoman, my daughter. And you, of singular clemency and pity --- when I had fallen, together with the commonwealth, by no tribunician but a consular blow --- were of such villainy and such intemperance that you would not allow even one hour to pass between my plague and your plunder, until at least the lamentation and the groaning of the city should fall silent!

\ciceroSection{18} It was not yet publicly known that the commonwealth had fallen, when the funeral arrangements were being settled to your account. At one and the same time my house was being plundered, was burning; my goods were being carried off from the Palatine to the consul next door, from my Tusculan villa to the other consul next door --- with the same hired ruffians acting as voters, with the same gladiator as carrier of the law, in a Forum empty not only of good men but even of free men, with the Roman people not knowing what was being done, the Senate suppressed and stricken: the treasury, provinces, legions, commands were being given as gifts to two impious and unspeakable consuls. The ruins these consuls left, you, the consuls, have shored up by your courage --- borne up by the great good faith and care of the tribunes of the plebs and of the praetors.

\ciceroSection{19} What shall I say of the most outstanding man, T. Annius? --- or who shall ever speak worthily enough of so great a citizen? Who, when he saw that a villainous citizen --- or rather a household enemy --- if it was permitted to use the laws, must be broken in court; but if violence itself impeded and removed the courts, then audacity must be overcome by courage, madness by fortitude, recklessness by judgement, brute force by resources, violence by violence --- first prosecuted him under the law on violence; then, when he saw that the courts had been removed by that same man, he took care that he might not be able to do everything by violence: who taught that neither houses nor temples nor the Forum nor the Curia could be defended without the highest courage and the greatest resources and forces from a domestic banditry; who first, after my withdrawal, drove fear from good men, hope from the audacious, terror from this order, and slavery from the citizenry.

\ciceroSection{20} Following this same plan with equal courage, spirit, and good faith, P. Sestius, for my safety, for your authority, for the standing of the state, never thought any enmities, any violence, any attacks, any peril of life had to be avoided. He so commended the cause of the Senate, harassed in the public meetings of villainous men, by his own care to the multitude, that nothing seemed so popular as your name, nothing once so dear to all as your authority. He defended me with everything that a tribune of the plebs could; the rest of his services he kept up just as if he had been my brother. By his clients, his freedmen, his household, his resources, his letters, I was so kept up that he seemed not only the helper of my calamity but its partner.
